% FRAGMENT A08 (drafter A: class model and main results)
% REPLACES: main.tex lines 566–686 ("\section{4. Equilibrium results}" to "… \hyperref[sec-coexistence-and-the-equilibrium-correspondence]{Section 4.3} describes the strengths between them.")
%   This covers the Section 4 heading and Section 4.1.
% LABELS KEPT: sec-results, sec-competition-creates-competition, result-proposition-2,
%   eq:paper-a2, eq:paper-a3, eq:paper-12.
% LABELS REMOVED: eq:paper-a1 (condition (A1) is deleted; the low cost is zero).
%   Old references to eq:paper-a1: main.tex 612, 1476, 1494, 1538 (all inside drafter A
%   fragments A08 and A13, where they are removed) and 1860 (inside fragment B02, which
%   replaces that line and no longer cites (A1)).
% LABELS ADDED: none.
% TAGS: (A2) and (A3) keep their tags and labels, so every \eqref elsewhere still works and
%   the hard-coded "(A2)", "(A3)" in the Online Appendix stay correct. Tag (A1) no longer
%   exists. If the author prefers no gap, retag (A2)->(A1) and (A3)->(A2) in one pass after
%   integration and update the hard-coded "(A2)", "(A3)" in the Online Appendix.
% NEW SYMBOLS: F_0, F_1, F_2 (incumbent laws of the three economies; uniform r_0<r_1<r_2 as
%   the leading case), e_H, e_L (class-conditional entry), W_H(F) (win probability of a
%   prepared class-H challenger), \mathsf O_H = \pi e_H W_H (high-class ownership).
% NEW CITATIONS: none.  NUM PLACEHOLDERS: none.
% SOURCES: Q1_N_main.md Claims 8, 10, 11; Q4_normalization.md Sections 3.6 and 7 (ratio form
%   of (A3)); SYNTHESIS.md A.1 rows 4, 5, 23.

\section{4. Equilibrium results}\label{sec-results}

\subsection{4.1 Competition creates
competition}\label{sec-competition-creates-competition}

The main result identifies an open set of economies in which stronger
competition changes the unique trading and on-path preparation outcome
from an uninformative price with low-cost preparation only to an
informative price with additional high-cost preparation.

\Needspace{8\baselineskip}

\protect\phantomsection\label{result-proposition-2}{\textbf{Proposition
2 (competition creates competition; analytical).}} \emph{Let the classes
satisfy the restrictions of
\hyperref[sec-values-preparation-and-the-sale]{Section 2.2}, let \(F_1\)
be a continuous first-order stochastic strengthening of the continuous
incumbent distribution \(F_0\) (for uniform incumbents, \(r_0<r_1\)),
fix \(0<\pi<1\), \(0<\rho<1\), \(b>1\), and \(k>0\), and suppose that}

\[
B_{F_0}(\pi)<c_H<B_{F_1}(M(\pi)),
\tag{A2}\label{eq:paper-a2}
\]

\[
\Delta_T(F_0)<k<\left(1-\frac1b\right)\rho\,\kappa(\pi)\,\Delta_T(F_1).
\tag{A3}\label{eq:paper-a3}
\]

\emph{(i) With the weak incumbent \(F_0\), the unique equilibrium
trading outcome is \(q_H=q_L=0\), the price carries no information, and
entry equals \(\rho\). (ii) With the strong incumbent \(F_1\), the
unique equilibrium trading outcome is \((q_H,q_L)=(1,-1)\), the price is
informative, and entry strictly exceeds \(\rho\). Let
\(W_H(F)=\mathbb E_H[F(\theta)\mathbf1\{\theta\ge p\}]\) be the
probability that a prepared class-\(H\) challenger wins, and let
\(e_H\) be entry conditional on class \(H\) at \(F_1\). High-class
ownership, the probability that a class-\(H\) challenger acquires the
target, rises from \(F_0\) to \(F_1\) if and only if
\(e_HW_H(F_1)>\rho W_H(F_0)\). It rises whenever
\((1+\rho)W_H(F_1)\ge2\rho W_H(F_0)>0\), and so in the point-mass
benchmark, where \(W_H=1\). (iii) For any continuous first-order
stochastic strengthening \(F_2\) of \(F_1\) with
\(B_{F_2}(M(\pi))<c_H\), the unique trading outcome is again
\((1,-1)\), but entry returns to \(\rho\). For every pair of classes
with \(\mathbb E_H[\max\{\theta,p\}]>\mathbb E_L[\max\{\theta,p\}]\)
and every \(\pi\), \(\rho\), and \(b\), these conditions, the condition
of part (iii), and the sufficient ownership condition hold together on a
nonempty open set of uniform strengths \(r_0<r_1<r_2\) and costs
\((c_H,k)\). Uniqueness refers to trading and on-path entry under
truthful bidding and allows arbitrary mixed orders and every unilateral
deviation to an order in \([-1,1]\).}

Because the weak-incumbent price is uninformative and the strong-incumbent price is informative, the strong-incumbent price experiment strictly Blackwell-dominates the weak one at the same noise law (\citeproc{ref-Blackwell1953}{Blackwell 1953}).

A low-cost challenger pays nothing to prepare, so it prepares after
every price. The proofs use this participation floor in three ways.
Entry is at least \(\rho\) at every price, so the denominator in
\eqref{eq:paper-10} is positive and the price reveals the market's
posterior. The investor's residual advantage is at least
\(\rho\kappa(\pi)\Delta_T\) at every flow, the bound in
\eqref{eq:paper-11} that \eqref{eq:paper-a3} compares with the trading
cost. Low-cost preparation does not respond to the price, so information
can only add entry.

The floor needs no assumption, because gross profit is an expected
auction rent and is never negative; any low cost below
\(B_F(m(\pi))\) in every economy compared would give the same orders,
prices, and entry.

Condition \eqref{eq:paper-a2} excludes high-cost preparation at the
weak-incumbent prior and permits it at favorable strong-incumbent
prices. Condition \eqref{eq:paper-a3} places the trading cost above
every informational return in the weak economy and below the global
bound on marginal trading profit in the strong economy. Its window is
nonempty exactly when
\(\Delta_T(F_0)/\Delta_T(F_1)<(1-1/b)\rho\kappa(\pi)\), a condition on
the two incumbent distributions, the classes, \(\rho\), \(b\), and
\(\pi\) that involves no unit of money. A prior away from one half
narrows the window, since \(\kappa(\pi)<m\) there.

These inequalities determine trading before imposing a particular order
profile. In the weak economy, the investor's gross advantage per unit is
at most \(\Delta_T(F_0)<k\), so every nonzero order loses money. With no
informed trading, the posterior stays at the prior and only low-cost
preparation occurs. In the strong economy, the residual lower bound is
large enough to make every increase in a correctly signed order
profitable. The bound controls the full order interval, including
deviations from mixed candidate strategies. Full correctly signed orders
are therefore necessary in every equilibrium.
\hyperref[paper-appendix]{Appendix A} derives the bound and constructs
the associated price and entry schedules.

Under full orders, favorable prices cross the expensive challenger's
preparation threshold with positive probability. With \(g_H\) and
\(g_L\) evaluated at \(F_1\), define the threshold belief \(\tau\), its
corresponding flow \(x^*\), the class-conditional probabilities of
crossing it \(\alpha_H,\alpha_L\), class-conditional entry \(e_H,e_L\),
total entry \(\mathsf E\), and high-class ownership \(\mathsf O_H\) by

\[
\begin{aligned}
\tau&=\frac{c_H-g_L}{g_H-g_L},&
x^*&=\frac b2\left[\log\frac\tau{1-\tau}-\log\frac\pi{1-\pi}\right],\\
\alpha_H&=1-\frac12e^{(x^*-1)/b},&
\alpha_L&=\frac12e^{-(x^*+1)/b},\\
e_H&=\rho+(1-\rho)\alpha_H,&
e_L&=\rho+(1-\rho)\alpha_L,\\
\mathsf E&=\pi e_H+(1-\pi)e_L,&
\mathsf O_H&=\pi e_HW_H(F_1).
\end{aligned}
\tag{12}\label{eq:paper-12}
\]

The prior shifts the threshold flow by \(-(b/2)\log[\pi/(1-\pi)]\). In
the benchmark that shift is zero and \(W_H=1\), because \(h\) exceeds
the incumbent's support, so
\(\mathsf E=\rho+(1-\rho)(\alpha_H+\alpha_L)/2\) and
\(\mathsf O_H=[\rho+(1-\rho)\alpha_H]/2\).

Throughout, entry \(\mathsf E\) is the preparation probability, the
probability that the challenger pays its preparation cost. When both
classes lie above the reserve, as in the benchmark, every prepared
challenger is an admissible bidder, and the auction has two admissible
bidders whenever the incumbent's value also meets the reserve.
\hyperref[sec-design]{Section 6.3} separates preparation from sale and
from two admissible bidders when the reserve is raised. Condition
\eqref{eq:paper-a2} places \(\tau\) in \((\pi,M(\pi))\) and \(x^*\) in
\((0,1)\). Since \(\alpha_H>1/2>\alpha_L\), the additional entry is
tilted toward the high class.

Ownership needs its own condition when the class values overlap the
incumbent's support. A stronger incumbent then beats a prepared
class-\(H\) challenger more often, which lowers \(W_H\), so entry can
rise while high-class ownership falls. The condition in part (ii)
compares the two effects. It holds whenever the high class rarely loses
to the incumbent, and automatically when its values lie above the
incumbent's support in both economies.

At a sufficiently strong incumbent satisfying part (iii), even the
largest feasible belief fails to justify high-cost preparation. Trading
remains informative, but entry returns to \(\rho\). The proposition establishes a rise and a subsequent fall across the specified economies; \hyperref[sec-coexistence-and-the-equilibrium-correspondence]{Section 4.3} describes the strengths between them.
